% ============================================================
% Testing for Drug-Enforcement Substitution After Marijuana
% Legalization: Marijuana vs. Alcohol Arrests
% Formatted to Economic Inquiry (Wiley/WEAI) submission requirements.
% Compile: pdfLaTeX -> BibTeX -> pdfLaTeX -> pdfLaTeX (Overleaf: Recompile)
% ============================================================
\documentclass[11pt]{article}

\usepackage[margin=1in]{geometry}
\usepackage{setspace}
\usepackage{amsmath,amssymb}
\usepackage{graphicx}
\usepackage{booktabs}
\usepackage{caption}
\usepackage[authoryear,round]{natbib}
\usepackage[hidelinks]{hyperref}
\usepackage{enumitem}
\usepackage{titlesec}

\doublespacing
\titleformat{\section}{\normalfont\bfseries\large}{\thesection.}{0.5em}{}
\titleformat{\subsection}{\normalfont\bfseries}{\thesubsection}{0.5em}{}

\title{}
\author{}
\date{}

\begin{document}

% ============================================================
% TITLE PAGE
% ============================================================
\begin{titlepage}
\begin{center}
\vspace*{1.5cm}
{\LARGE \bfseries Testing for Drug-Enforcement Substitution After\\
Marijuana Legalization\par}
\vspace{0.3cm}
{\large \itshape Marijuana vs. Alcohol Arrest Trends, 2005--2023\par}
\vspace{1.2cm}
{\large Author name(s) withheld for review\par}
\vspace{0.3cm}
{\normalsize Working Paper --- July 2026\par}
\end{center}

\vspace{1.2cm}
\noindent\textbf{Abstract} \\[4pt]
\noindent We test for substitution between marijuana and alcohol
using FBI arrest data: marijuana (possession/sale) and
alcohol-specific (liquor-law/drunkenness) offenses. We use (i) a
switcher/control county design in never-legalizing states, with
treatment defined as the raw change in miles to the nearest legal
border, and (ii) a 21-state staggered-adoption panel around each
state's own legalization date. We find no evidence of substitution: a
100-mile reduction in border distance leaves marijuana arrests
unchanged (p\,=\,0.94) but significantly increases alcohol arrests
(p\,=\,0.04) -- the opposite sign from substitution. The state panel
shows both series on a common secular decline, undisturbed by
legalization timing (p\,=\,0.95, 0.47). \\[10pt]
\noindent\textbf{JEL Classification:} K42, I18, R12 \\
\noindent\textbf{Keywords:} recreational marijuana; alcohol; drug
substitution; arrest data; NIBRS; difference-in-differences

\vfill
\noindent\textbf{Abbreviations} \\[2pt]
\noindent
\begin{tabular}{@{}ll@{}}
DiD    & Difference-in-differences \\
FBI    & Federal Bureau of Investigation \\
NIBRS  & National Incident-Based Reporting System \\
SE     & Standard error \\
TWFE   & Two-way fixed effects \\
\end{tabular}

\vspace{0.5cm}
\noindent{\footnotesize A separate Disclosure Statement accompanies this
submission per Economic Inquiry policy and is not reproduced in this file.}
\end{titlepage}

% ============================================================
% 1. INTRODUCTION
% ============================================================
\section{Introduction}

A long-standing question in drug policy is whether marijuana and
alcohol are economic substitutes or complements. If substitutes, legal
marijuana access should reduce alcohol consumption and, mechanically,
alcohol-related law enforcement contact (drunkenness and liquor-law
arrests). \citet{AndersonHansenRees2013} find evidence consistent with
substitution -- medical marijuana laws are associated with reduced
traffic fatalities and lower alcohol consumption, particularly among
young adults. This paper asks whether the same substitution pattern is
visible in \emph{arrest} data around the 2020--2023 recreational
legalization wave, using the same never-legalizing-state panel and
continuous distance-to-border design developed in the companion
project on marijuana-related driving spillovers
\citep{HaoCowan2020,Moyer2026}.

We pursue two complementary designs. First, within never-legalizing
states, we compare marijuana- and alcohol-arrest trends in counties
that became closer to a newly-legal border (``switchers'') against
counties that did not, exactly as in the companion driving-spillover
analysis but with arrest outcomes rather than fatal crashes. Second,
we examine the 21 states that themselves legalized recreational
marijuana between 2014 and 2024, comparing their own marijuana and
alcohol arrest trends before and after their individual legalization
dates in a staggered-adoption panel.

We find no robust evidence of substitution in either design. If
anything, the county-level continuous-distance specification finds a
statistically significant \emph{increase} in alcohol arrests among
counties that became closer to a legal border -- the opposite sign
from what a substitution story predicts -- while the marijuana-arrest
response is negative but not significant. The state-level
staggered-adoption panel finds both series on a common declining
trend with no detectable break at each state's own legalization date.
We also document, and correct for, a serious data-quality issue: the
FBI's mandatory 2021 transition to NIBRS-only reporting caused a sharp,
temporary collapse in state-level reporting coverage (to as low as 36\%
in one state) that can mechanically distort any raw pre/post arrest
comparison spanning that year.

% ============================================================
% 2. DATA
% ============================================================
\section{Data}

\subsection{Arrest Measures}

We use the FBI's Crime Data Explorer (CDE) API, which reports monthly
arrest counts by specific offense code at the national, state, and
agency level. \textbf{Marijuana arrests} combine Drug Possession,
Marijuana (code 158) and Drug Sale/Manufacturing, Marijuana (code
153). \textbf{Alcohol arrests} combine Liquor Law Violations (code
270) and Drunkenness (code 280) -- offenses that, unlike Driving Under
the Influence (code 260), are attributable to a specific substance
rather than pooling alcohol- and drug-impaired driving together. This
mirrors the marijuana-code construction and avoids the substance-
attribution problem documented in the companion driving-spillover
paper.

\subsection{A Critical Data-Quality Issue: the 2021 NIBRS Transition}

The FBI required all reporting agencies to transition from the legacy
Summary Reporting System to the National Incident-Based Reporting
System effective January 2021. Many agencies were not yet NIBRS-
certified and temporarily stopped reporting. We document this directly:
national population coverage fell from 96.7\% (January 2020) to 74.6\%
(January 2021); one state in our legalizing-state panel (Illinois) saw
coverage collapse from 94.1\% to 36.4\% over the same window, before
recovering to 90.9\% by December 2023. Because the composition of
agencies that continue reporting during a coverage collapse is not
random -- larger, better-resourced (often more urban, higher-
enforcement-intensity) agencies are more likely to remain NIBRS-
compliant -- raw arrest rates during the transition window are not
comparable to rates before or after it, even after adjusting for
reporting population. Figure~\ref{fig:co} illustrates this for Colorado.

We address this in two ways. First, our primary switcher/control
comparison uses only 2020 and 2023 -- both outside the acute transition
window -- and we confirmed that 2020-to-2023 coverage change is small
(under 5 percentage points) for all 21 legalizing states and all but a
handful of the 32 never-legalizing states (Florida and Pennsylvania
show the largest shifts, $-19.3$ and $+14.9$ points respectively;
results are robust to excluding them, available on request). Second,
our state-panel event-study design includes year fixed effects, which
absorb the common national 2021 coverage shock, identifying the
legalization effect only from the \emph{relative} timing of
legalization across states.

\begin{figure}[htbp]
  \centering
  \includegraphics[width=0.85\textwidth]{figures/figB2_colorado.png}
  \caption{Colorado Monthly Marijuana and Alcohol Arrest Rates, 2005--2023. The shaded region marks the 2021--2022 NIBRS transition, during which state reporting coverage fell to roughly 54\%.}
  \label{fig:co}
\end{figure}

% ============================================================
% 3. DESIGN A: SWITCHER/CONTROL COUNTIES
% ============================================================
\section{Design A: Switcher vs. Control Counties}

We reuse the companion project's never-legalizing-state panel: 32
states that had not legalized recreational marijuana through 2023,
with each county's distance to the nearest legal-market border
computed in miles for 2020 and 2023 (Haversine distance between
county centroids). Agency-level arrest counts were pulled from the
FBI CDE API for all agencies in these states (13,384 agencies),
matched to Census county FIPS by normalized county name (83.5\% match
rate among single-county agencies), and aggregated to the county
level. This yields matched arrest data for 1,793 counties in the main
sample (378 switchers, 1,415 controls) -- 82.8\% of the
never-legalizing-state county panel.

\subsection{The Continuous Treatment: Miles to the Nearest Border}

Our primary treatment variable is the county's raw change in distance
to the nearest recreational-marijuana state border, in miles:
\begin{equation}
  \Delta\text{Dist}_i \;=\; \text{Dist}_{i,2023} \;-\; \text{Dist}_{i,2020},
  \label{eq:distmi}
\end{equation}
which is negative for counties that became closer to a newly-legal
border and positive for counties that fell further behind (e.g.,
because a once-nearby state never opened retail sales while others
did). Figure~\ref{fig:map} maps the resulting treatment assignment,
and Figure~\ref{fig:distdist} shows the distribution of both the level
(2023 distance) and the change ($\Delta\text{Dist}_i$) separately for
switchers and controls -- switchers are, by construction, the
counties whose distance fell below the 100-mile threshold and whose
$\Delta\text{Dist}_i$ is concentrated in the $-400$ to $0$ mile range.

\begin{figure}[htbp]
  \centering
  \includegraphics[width=0.85\textwidth]{figures/figA0a_map.png}
  \caption{Treatment Assignment: Switcher and Control Counties by Location.}
  \label{fig:map}
\end{figure}

\begin{figure}[htbp]
  \centering
  \includegraphics[width=0.95\textwidth]{figures/figA0b_distance_dist.png}
  \caption{The Continuous Treatment: Distribution of Distance (Level and Change) to the Nearest Recreational Border.}
  \label{fig:distdist}
\end{figure}

\begin{table}[htbp]
\centering
\caption{Summary Statistics}
\begin{tabular}{@{}lc@{}}
\toprule
Variable & Value \\
\midrule
Marijuana arrest rate per 100k (2020) & 10.74 \\
Marijuana arrest rate per 100k (2023) & 7.78 \\
Alcohol arrest rate per 100k (2020) & 42.64 \\
Alcohol arrest rate per 100k (2023) & 39.96 \\
Switcher counties & 378 \\
Control counties & 1{,}415 \\
\bottomrule
\end{tabular}
\end{table}

\begin{figure}[htbp]
  \centering
  \includegraphics[width=0.95\textwidth]{figures/figA1_trends.png}
  \caption{Marijuana vs. Alcohol Arrest Rate Trends, Switchers vs. Controls.}
  \label{fig:trendsA}
\end{figure}

Table~\ref{tab:regA} reports first-difference regressions (algebraically
equivalent to county $+$ year fixed effects for $T=2$):
\begin{equation}
  \Delta Y_i \;=\; \alpha \;+\; \beta \cdot \frac{\Delta\text{Dist}_i}{100} \;+\; \varepsilon_i,
  \label{eq:distreg}
\end{equation}
scaling $\Delta\text{Dist}_i$ into 100-mile units so that $\beta$ is
directly readable as ``change in arrest rate per 100,000 associated
with a 100-mile change in distance to the nearest border,'' alongside
the binary switcher indicator and, as a robustness check, the
$\Delta\log(\text{dist})$ specification used in the companion
driving-spillover paper. SE clustered by state throughout.

\begin{table}[htbp]
\centering
\caption{Effect on Marijuana and Alcohol Arrest Rates}
\label{tab:regA}
\begin{tabular}{@{}lcc@{}}
\toprule
 & Coefficient (SE) & $p$-value \\
\midrule
\textit{Marijuana} & & \\
\quad Switcher (binary) & $-3.177$ (2.308) & 0.169 \\
\quad Distance, per 100 mi [\textbf{primary}] & $-0.064$ (0.798) & 0.936 \\
\quad Log-distance (robustness) & $+1.161$ (1.540) & 0.451 \\
\textit{Alcohol} & & \\
\quad Switcher (binary) & $+4.845$ (5.894) & 0.411 \\
\quad Distance, per 100 mi [\textbf{primary}] & $-5.953^{**}$ (2.914) & 0.041 \\
\quad Log-distance (robustness) & $-6.836^{**}$ (2.669) & 0.010 \\
\bottomrule
\end{tabular}
\vspace{2pt}
\begin{minipage}{0.9\textwidth}
\footnotesize Note: Outcome is the change in arrest rate per 100{,}000,
2020--2023. ``Distance, per 100 mi'' is $\beta$ from
Equation~\ref{eq:distreg}: a negative $\beta$ means the arrest rate
\emph{fell} as distance to the border \emph{increased}, i.e., the rate
\emph{rose} as the county became closer. SE clustered by state
($N=1{,}793$). $^{**}p<0.05$.
\end{minipage}
\end{table}

The primary, directly-interpretable miles specification shows no
detectable relationship for marijuana arrests: a 100-mile change in
distance to the border is associated with a $-0.064$ change in the
arrest rate per 100{,}000, statistically indistinguishable from zero
($p=0.936$). For alcohol, the coefficient is $-5.953$ and significant
at the 5\% level ($p=0.041$): because the sign is negative on
\emph{distance} (not distance-\emph{closeness}), this means the
alcohol arrest rate \emph{rose} as counties became 100 miles closer to
a newly-legal border. This is the opposite sign from a substitution
story, which would predict alcohol arrests to \emph{fall} as
marijuana becomes more accessible. The log-distance robustness check
(bottom row of each panel) agrees in sign and significance with the
miles-based primary specification for both outcomes. The binary
switcher specification points in the same qualitative direction for
both outcomes (marijuana $-3.177$, alcohol $+4.845$) but neither
reaches significance -- unsurprising, since the binary indicator
discards the within-group variation in exactly how many miles closer
each switcher became, which Figure~\ref{fig:distdist}(b) shows is
substantial.

\begin{figure}[htbp]
  \centering
  \includegraphics[width=0.95\textwidth]{figures/figA4_treatment_scatter.png}
  \caption{The Continuous-Treatment Regression Itself: Distance Change vs. Arrest-Rate Change.}
  \label{fig:treatscatter}
\end{figure}

Figure~\ref{fig:treatscatter} shows the regression underlying
Table~\ref{tab:regA}'s primary specification directly: each point is a
county, with its change in distance to the border (miles) on the
horizontal axis and its change in arrest rate on the vertical axis.
The near-flat fitted line for marijuana and the modestly downward-
sloping line for alcohol (moving left, i.e.\ toward more negative
$\Delta\text{Dist}$, the fitted alcohol line rises) are the visual
counterpart of the coefficients in Table~\ref{tab:regA}.

\begin{figure}[htbp]
  \centering
  \includegraphics[width=0.7\textwidth]{figures/figA2_coefplot.png}
  \caption{Switcher Effect on Arrest Rates with 95\% Confidence Intervals.}
  \label{fig:coefA}
\end{figure}

\begin{figure}[htbp]
  \centering
  \includegraphics[width=0.65\textwidth]{figures/figA3_scatter.png}
  \caption{County-Level Co-Movement of Marijuana and Alcohol Arrest-Rate Changes.}
  \label{fig:scatterA}
\end{figure}

As a further check, we compute the raw correlation between
county-level changes in the two arrest rates: $r=0.12$, a weak
\emph{positive} co-movement (Figure~\ref{fig:scatterA}). A clean
substitution story predicts a negative correlation (counties where
marijuana arrests fall most should see alcohol arrests rise most); a
weak positive correlation is more consistent with both series
responding to a common local factor (e.g., overall policing intensity
or economic activity near the border) than with behavioral
substitution between the two goods.

\subsection{Heterogeneity by Age of Offender}
\label{sec:agecohort}

The county-level result above -- alcohol arrests rising as counties
grow closer to a legal marijuana border -- is the paper's only
statistically significant finding, and it points toward complementarity
rather than substitution. A natural follow-up question is whether this
pattern is concentrated in a particular age cohort: young adults newly
able to buy marijuana legally, for instance, might drive the effect,
or it could be spread evenly across ages (or driven by older
offenders instead). We re-estimate Design A's switcher/continuous-
distance specification separately by age of arrestee, holding the
treatment design fixed.

\subsubsection{Data: NIBRS Arrestee-Level Microdata}

Age of offender is not available through the FBI Crime Data Explorer's
live query API used elsewhere in this paper -- that API returns
offense-level monthly counts only, with no demographic breakdown. We
instead use the ICPSR/NACJD ``NIBRS Extract Files'' (2020: ICPSR 38566
v3; 2023: ICPSR 39270 v2), a cleaned, arrestee-level repackaging of the
raw National Incident-Based Reporting System microdata in which each
row is one arrestee, carrying the linked incident's offense and
property detail alongside the arrestee's own age. Marijuana arrests
are identified as UCR arrest offense code 351 (Drug/Narcotic
Violations) with a linked ``suspected drug type'' of marijuana
recorded anywhere on the incident's property segment; alcohol arrests
are UCR arrest offense codes 905 (Drunkenness) and 907 (Liquor Law
Violations), again excluding Driving Under the Influence (904) for the
same substance-attribution reason as the main analysis. We restrict to
single-county agencies in the 32 never-legalizing states, matching
Design A, and further restrict to the 1,861 counties with an active,
single-county, NIBRS-reporting agency in \emph{both} 2020 and 2023, so
that a county with zero recorded marijuana or alcohol arrests in a
given age bin can be read as a genuine zero rather than simply missing
NIBRS coverage. Age is grouped into four bins: under 21 (below the
legal alcohol-purchase age), 21--30, 31--45, and 46 and older.

\subsubsection{A Data-Quality Caveat: NIBRS Microdata vs.\ the SRS-Style API}

Restricted to the identical set of counties and year, the NIBRS
microdata used in this section and the FBI CDE API counts used
everywhere else in the paper disagree substantially in scale --
NIBRS yields roughly 13 times more marijuana arrests and roughly 9
times fewer alcohol arrests. We could not fully reconcile this from
public documentation, but two mechanisms plausibly explain most of
it. First, the FBI's classic SRS-style arrest classification is
believed to apply a substance-severity hierarchy when an arrest
incident lists more than one drug (e.g., marijuana recovered alongside
cocaine is classified under the more serious substance), which we
cannot replicate from raw NIBRS incident-level drug-type flags alone;
this would mechanically deflate SRS-style marijuana counts relative to
a full incident-level accounting like the one used here. Second, NIBRS
Group~B arrest-only offenses -- which include both Drunkenness and
Liquor Law Violations -- are documented to have materially weaker
agency reporting compliance than incident-linked Group~A offenses,
consistent with the undercount we observe. Because of this, the
age-cohort results below should be read as an internally consistent
comparison \emph{across} age bins within the NIBRS source, not as a
like-for-like replication of the paper's SRS-based headline
coefficients in Table~\ref{tab:regA}. We also do not observe
age-specific county population, so all rates use total county
population as the denominator throughout; this affects the
\emph{level} of each bin's rate but not the within-bin switcher/control
comparison, which is what the regression coefficients capture.

\subsubsection{Results}

\begin{table}[htbp]
\centering
\caption{Effect on Marijuana and Alcohol Arrest Rates, by Age of Arrestee}
\label{tab:regC}
\begin{tabular}{@{}lcccc@{}}
\toprule
Age bin & \multicolumn{2}{c}{Marijuana} & \multicolumn{2}{c}{Alcohol} \\
 & Switcher & Dist., per 100 mi & Switcher & Dist., per 100 mi \\
\midrule
Under 21   & $-5.580$ (5.287)   & $+0.554$ (1.872)  & $+0.146$ (0.476)  & $+0.155$ (0.193) \\
21--30     & $-15.129$ (11.601) & $+2.121$ (4.000)  & $-0.235$ (0.202)  & $+0.067$ (0.071) \\
31--45     & $-7.700$ (7.669)   & $+0.310$ (2.806)  & $-0.240$ (0.204)  & $+0.014$ (0.081) \\
46 and older & $-0.280$ (3.125) & $-0.213$ (1.069)  & $-0.057$ (0.176)  & $-0.034$ (0.045) \\
\bottomrule
\end{tabular}
\vspace{2pt}
\begin{minipage}{0.95\textwidth}
\footnotesize Note: Cell entries are coefficient (clustered SE); none
are statistically significant at conventional levels (all $p>0.19$).
$N=1{,}736$ counties per cell (399 switchers, 1{,}337 controls) in
every age bin -- the county set is fixed across bins by construction
(Section~3.4.1), so only the arrest counts vary. Outcome, sample, and
clustering match Table~\ref{tab:regA}.
\end{minipage}
\end{table}

\begin{figure}[htbp]
  \centering
  \includegraphics[width=0.8\textwidth]{figures/figC1_age_coefplot.png}
  \caption{Primary Continuous-Distance Coefficient by Age of Arrestee, with 95\% Confidence Intervals.}
  \label{fig:coefC}
\end{figure}

No coefficient in Table~\ref{tab:regC} is statistically distinguishable
from zero, for either outcome, in any age bin. Most relevant to the
paper's headline result: the alcohol distance coefficient -- significant
and positive-on-closeness in the SRS-based Design A analysis
(Table~\ref{tab:regA}) -- is small and statistically null in every age
bin here, ranging from $-0.034$ to $+0.155$ with no consistent sign
across bins (Figure~\ref{fig:coefC}). Given the data-quality caveat
above, we do not read this as evidence that the complementarity signal
is concentrated in, or absent from, any particular age cohort. The
underlying alcohol cell counts are thin -- total NIBRS-recorded
alcohol arrests per age bin across all 1,736 counties and both years
range from roughly 200 to 1,900 -- and a measure this sparse is likely
underpowered to detect age heterogeneity in an effect that was only
marginally significant ($p=0.041$) in the much larger SRS-based sample
to begin with. We view the age-disaggregated NIBRS alcohol measure, as
constructed here, as insufficient to answer the age-cohort question
either way; a more complete answer would require either a less sparse
alcohol-enforcement measure or the official NIBRS-to-SRS drug
classification crosswalk needed to reconcile the scale discrepancy
noted above.

% ============================================================
% 4. DESIGN B: STAGGERED ADOPTION ACROSS LEGALIZING STATES
% ============================================================
\section{Design B: Legalizing States' Own Before/After}

For the 21 states that legalized recreational marijuana between 2014
and 2024, we pulled state-level monthly arrest data from January 2005
through December 2023 (228 months per state) and constructed an
annual panel (21 states $\times$ 19 years $=$ 399 observations). We
estimate
\begin{equation}
  y_{it} = \alpha_i + \gamma_t + \beta \cdot \text{Post}_{it} + \varepsilon_{it},
  \label{eq:twfe}
\end{equation}
with state fixed effects $\alpha_i$, year fixed effects $\gamma_t$
(which absorb the common national 2021 reporting shock), and
$\text{Post}_{it}=1$ in and after each state's own legalization year,
SE clustered by state.

\begin{table}[htbp]
\centering
\caption{Staggered-Adoption TWFE: State's Own Legalization}
\label{tab:twfeB}
\begin{tabular}{@{}lccc@{}}
\toprule
Outcome & $\beta$ (Post) & SE & $p$-value \\
\midrule
Marijuana arrest rate & $-0.029$ & 0.426 & 0.946 \\
Alcohol arrest rate & $-1.950$ & 2.693 & 0.469 \\
\bottomrule
\end{tabular}
\vspace{2pt}
\begin{minipage}{0.8\textwidth}
\footnotesize Note: $N=399$ (21 states $\times$ 19 years). State and
year fixed effects. SE clustered by state.
\end{minipage}
\end{table}

Neither coefficient is distinguishable from zero. This may seem to
contradict the large raw pre/post gap (marijuana: 1.70 $\to$ 0.37 per
100{,}000; alcohol: 12.03 $\to$ 4.08), but Figure~\ref{fig:eventB}
resolves the apparent tension: pooling all 21 states by
years-since-own-legalization, both series decline \emph{smoothly and
continuously} for years before legalization and show no discontinuity
at event-year zero. The raw pre/post gap reflects a common secular
decline in both offense categories across essentially all states
(legalizing and not) over 2005--2023 -- plausibly reflecting broader
shifts in enforcement priorities and criminal-justice reform -- not an
effect specific to each state's own legalization timing.

\begin{figure}[htbp]
  \centering
  \includegraphics[width=0.95\textwidth]{figures/figB1_eventstudy.png}
  \caption{Event Study: Arrest Rates Around Each State's Own Legalization Date.}
  \label{fig:eventB}
\end{figure}

% ============================================================
% 5. DISCUSSION
% ============================================================
\section{Discussion}

Across two independent designs -- a within-state continuous-distance
comparison and a between-state staggered-adoption panel -- we find no
consistent evidence that marijuana legalization substitutes away from
alcohol as measured by arrest data. The one statistically significant
result we obtain (Design A, continuous distance, alcohol arrests) runs
in the direction opposite to substitution. We view this as suggestive
rather than conclusive: arrest counts reflect both underlying behavior
and enforcement intensity, and a rise in alcohol arrests near a new
legal-market border could reflect increased cross-border commercial
and social activity (more visitors, more nightlife, more overall
police presence) rather than a behavioral shift toward heavier
drinking. This is analogous to the concern raised in
\citet{HaoCowan2020} that arrest-based measures conflate enforcement
effort with underlying behavior.

The state-panel null result (Design B) is, in our view, the more
credible of the two: it directly controls for the strong common
secular decline in both offense categories via year fixed effects,
whereas the raw pre/post comparison that motivates a naive
substitution narrative is confounded by that same secular trend. We
therefore conclude that arrest data, at least as constructed here, do
not support a robust marijuana-alcohol substitution story around the
2014--2024 legalization wave -- a useful complement to
\citet{AndersonHansenRees2013}'s consumption-based evidence, since
arrest data capture a different (enforcement-mediated) margin of the
same underlying question.

% ============================================================
% 6. LIMITATIONS
% ============================================================
\section{Limitations}

\begin{enumerate}[label=(\arabic*), leftmargin=*]
\item \textbf{Arrests reflect enforcement, not just behavior.} A
change in arrest rates conflates true behavioral change with changes
in policing priorities, which may themselves respond to legalization
(e.g., reduced marijuana enforcement effort could mechanically free up
capacity for alcohol enforcement, independent of any behavioral
substitution).
\item \textbf{The 2021 NIBRS transition is a first-order confound}
for any raw pre/post design spanning that year; we mitigate but cannot
fully eliminate this concern, particularly for states or counties with
large 2020-to-2023 coverage changes (Florida, Pennsylvania).
\item \textbf{County-level agency matching is incomplete.} Only 83.5\%
of single-county agencies matched to a Census FIPS code by name;
multi-county agencies (e.g., state police) were excluded entirely from
the county-level rollup.
\item \textbf{Liquor-law and drunkenness arrests are an imperfect proxy
for alcohol-related harm.} These offenses are sensitive to local
enforcement norms (e.g., open-container and public-intoxication
policing varies widely by jurisdiction) independent of actual alcohol
consumption.
\item \textbf{Two-period main design.} As in the companion project,
the switcher/control comparison uses only 2020 and 2023, precluding a
formal pre-trend test beyond the coverage-stability check reported in
Section 2.2.
\item \textbf{The age-cohort extension uses a different, not fully
reconciled, data source.} Section~3.4's age-of-offender breakdown uses
raw NIBRS arrestee microdata rather than the SRS-style API used
elsewhere, and the two disagree in scale by roughly an order of
magnitude for both outcomes (Section~3.4.2). The age-disaggregated
results should be read as internally consistent across age bins, not
as a decomposition of the paper's headline coefficients.
\end{enumerate}

% ============================================================
% 7. CONCLUSION
% ============================================================
\section{Conclusion}

We find no robust evidence that recreational marijuana legalization
induces substitution away from alcohol as measured by arrest data,
across two complementary quasi-experimental designs. Where a
statistically significant effect appears, it runs opposite to the
substitution hypothesis. Combined with the well-known 2021 NIBRS
reporting disruption we document here, we recommend that any future
work in this space (a) use state and year fixed effects to absorb
common reporting-coverage shocks, and (b) treat arrest-based
substitution evidence as complementary to, rather than a replacement
for, consumption- and health-outcome-based measures such as those in
\citet{AndersonHansenRees2013}.

% ============================================================
% REFERENCES
% ============================================================
\bibliographystyle{apalike}
\bibliography{references}

\end{document}
